% Generated by experiments/scripts/export_latex.py (export_latex_c12.py) -- do not edit by hand.
% Source: experiments/data/2-processed-tables-and-evaluations/revised-contract-12m/window-w1/eval_summary.json
% Source: experiments/data/2-processed-tables-and-evaluations/revised-contract-12m/supplementary/supplementary.json
\begin{table}[htbp]
\centering
\caption[기간 W1에서 엔도스랭크와 활동 재시작을 쓰는 그 변형.]{기간 W1의 스펜더 코호트와 라벨에서 엔도스랭크와, 활동 재시작을 쓰는 엔도스랭크 변형을 보여 줘요(계약 $n=7{,}915$개). 이 행들도, 그 아래의 대비들도 규칙 B의 적용한 두 기준에 들지 않아요. 그리고 이 대비들은 등록 문서에 이름이 적혀 있지 않았어요. 활동 재시작을 쓰는 엔도스랭크의 대비는 저장해 둔 W1 점수로 나중에 계산했어요. AWP, C-PR, 차수 행은 표~\ref{tab:fresh-holdout}의 값을 다시 적은 것이고, 모든 대비는 같은 계약 다시 뽑은 표본을 써요. 엔도스랭크와 AWP 사이의 차이는 표~\ref{tab:endorserank-awp-holdout}에, C-PR의 차이는 표~\ref{tab:cpr-contrasts}에 있어요. 신뢰구간(CI)은 계약을 다시 뽑아 구한 95\% 백분위 부트스트랩 구간이에요(짝지은 다시 뽑기 400번, 시드 42).}
\label{tab:fresh-posthoc}
\fitwidth{%
\begin{tabular}{lcc}
\toprule
2026년 6월 30일의 점수 & 새 허락 쌍 ($\tau$ [95\% 신뢰구간]) & 새 송금자 ($\tau$ [95\% 신뢰구간]) \\
\midrule
동결일의 받은 허락 수 (기준선) & 0.181 [0.158, 0.208] & 0.072 [0.046, 0.096] \\
동결일의 송금 입차수 (기준선) & 0.031 [0.008, 0.052] & 0.302 [0.278, 0.323] \\
\midrule
EndorseRank & 0.147 [0.125, 0.170] & 0.064 [0.041, 0.085] \\
EndorseRank, 활동 재시작(허락하는 활동) & 0.148 [0.126, 0.171] & 0.072 [0.049, 0.093] \\
AWP & 0.064 [0.042, 0.084] & 0.314 [0.295, 0.332] \\
C-PR ($\lambda=0.5$) & 0.104 [0.083, 0.126] & 0.236 [0.217, 0.255] \\
\midrule
EndorseRank 빼기 받은 허락 수, 새 허락 쌍 & \multicolumn{2}{l}{$\Delta\tau=-0.034$ [-0.043, -0.025]} \\
AWP 빼기 송금 입차수, 새 송금자 & \multicolumn{2}{l}{$\Delta\tau=+0.011$ [0.001, 0.022]} \\
\midrule
EndorseRank (활동 재시작) 빼기 받은 허락 수, 새 허락 쌍 & \multicolumn{2}{l}{$\Delta\tau=-0.033$ [-0.044, -0.022]} \\
EndorseRank (활동 재시작) 빼기 EndorseRank, 새 허락 쌍 & \multicolumn{2}{l}{$\Delta\tau=+0.001$ [-0.002, 0.005]} \\
EndorseRank (활동 재시작) 빼기 EndorseRank, 새 송금자 & \multicolumn{2}{l}{$\Delta\tau=+0.008$ [0.004, 0.012]} \\
\bottomrule
\end{tabular}
}
\end{table}
